\documentclass[
	degree          = bachelor,
	StudentName     = 张三,
	StudentID       = 2020123456,
	AdvisorName     = 李四,
	Grade           = 2020,
	Major           = 计算机科学与技术,
	Department      = 计算机学院,
	SubmitYear		= 2024,
	SubmitMonth		= 6,
	Title           = 基于XXX的YYY研究,
	TitleEng        = {{Research on YYY Based on XXX}},  % 选填，显示在第二封面
]{cauc_thesis}

% 也可以使用 \caucsetup 在导言区进行配置（与 \documentclass 选项等效）：
% \caucsetup{
%   Type=design,          % 本科毕业设计（默认 thesis）
%   declaration=true,     % 包含原创性声明
%   draft=false,          % 草稿模式
% }

\begin{document}

\frontmatter

\makecover

\begin{abstract-zh}
	摘要应具有独立性和自含性，即不需阅读报告、论文的全文，就能获得必要的信息。摘要是一篇完整的短文，一般应说明研究工作目的、实验方法、结果和最终结论等，而重点是结果和结论。

	摘要要求结构严谨，表达简明，语义确切，字数300-500字为宜。

	\keywordszh{关键词1；关键词2；关键词3；关键词4}
\end{abstract-zh}

\begin{abstract-en}
	English abstract content goes here. It should correspond to the Chinese abstract above.

	\keywordsen{Keyword1; Keyword2; Keyword3; Keyword4}
\end{abstract-en}

\makefrontmatter

\mainmatter
\pagestyle{fancy}

\chapter{绪论}

\section{研究背景}

研究背景内容。

\section{研究目的与意义}

研究目的与意义内容。

\chapter{结论}

结论内容。

\appendix
\chapter{附录标题示例}

这是附录的内容。附录章节编号使用大写字母（附录A、附录B……）。

\begin{table}[!ht]
	\centering
	\caption{附录表示例}\label{tab:appendix}
	\begin{tabular}{cc}
		\toprule 项目 & 结果 \\ \midrule
		数据1 & 0.95 \\
		数据2 & 0.87 \\
		\bottomrule
	\end{tabular}
\end{table}

\chapter{核心代码}

\begin{lstlisting}[language=Python, caption={二分查找实现}]
def binary_search(arr, target):
    left, right = 0, len(arr) - 1
    while left <= right:
        mid = (left + right) // 2
        if arr[mid] == target:
            return mid
        elif arr[mid] < target:
            left = mid + 1
        else:
            right = mid - 1
    return -1
\end{lstlisting}

\backmatter

\bibliography{reference}

\begin{acknowledgment}
	致谢内容。
\end{acknowledgment}

\end{document}
